\section*{习题 7.1\quad 微分方程的基本概念}

1. 指出下列各微分方程的阶数，并判断它们是否为线性方程

1. 题目：\(x^{2}y''-xy'+3y=\cos x\)

解答：
\begin{itemize}
\item 阶数：最高阶导数为 \(y''\)，故为二阶微分方程。
\item 线性性：方程中 \(y''\)、\(y'\) 和 \(y\) 都是一次项，且系数是 \(x\) 的函数，没有出现 \(y\) 或其导数的乘积或非线性函数，因此是线性方程。
\end{itemize}
答案：二阶线性方程

2. 题目：\(x^{2}dx=y^{3}dy\)

解答：
\begin{itemize}
\item 阶数：方程可改写为 \(\frac{dy}{dx}=\frac{x^{2}}{y^{3}}\)，最高阶导数为 \(\frac{dy}{dx}\)，故为一阶微分方程。
\item 线性性：方程中含有 \(y^{3}\)，是 \(y\) 的非线性项，因此是非线性方程。
\item 定义域说明：当 \(y=0\) 时方程无意义，应在 \(y\ne0\) 的区间内考虑。
\end{itemize}
答案：一阶非线性方程

3. 题目：\((1+y^{2})y'''+6(y'')^{2}+3y=0\)

解答：
\begin{itemize}
\item 阶数：最高阶导数为 \(y'''\)，故为三阶微分方程。
\item 线性性：方程中含有 \((y'')^{2}\)，是 \(y''\) 的非线性项，因此是非线性方程。
\end{itemize}
答案：三阶非线性方程

4. 题目：\(y''+\sin(x+y)=\sin x\)

解答：
\begin{itemize}
\item 阶数：最高阶导数为 \(y''\)，故为二阶微分方程。
\item 线性性：方程中含有 \(\sin(x+y)\)，是 \(y\) 的非线性函数，因此是非线性方程。
\end{itemize}
答案：二阶非线性方程

5. 题目：\(y^{(m)}+y''+y=0\ (m>2)\)

解答：
\begin{itemize}
\item 阶数：最高阶导数为 \(y^{(m)}\)，故为 \(m\) 阶微分方程。
\item 线性性：方程中 \(y^{(m)}\)、\(y''\) 和 \(y\) 都是一次项，且系数为常数，没有非线性项，因此是线性方程。
\end{itemize}
答案：\(m\) 阶线性方程

6. 题目：\(y''+p(x)y'+q(x)y=g(x)\)

解答：
\begin{itemize}
\item 阶数：最高阶导数为 \(y''\)，故为二阶微分方程。
\item 线性性：方程中 \(y''\)、\(y'\) 和 \(y\) 都是一次项，且系数是 \(x\) 的函数，没有非线性项，因此是线性方程。
\end{itemize}
答案：二阶线性方程

2. 指出下列各题中的函数是否为对应微分方程的解

1. 题目：\(xy'=x^{2}+y^{2}+y,\ y=x\tan\left(x+\frac{\pi}{6}\right)\)

解答：需要验证函数是否满足微分方程。

计算导数：
\[ y'=\tan\left(x+\frac{\pi}{6}\right)+x\sec^{2}\left(x+\frac{\pi}{6}\right). \]
代入方程左边：
\[ xy'=x\left[\tan\left(x+\frac{\pi}{6}\right)+x\sec^{2}\left(x+\frac{\pi}{6}\right)\right]. \]
右边：
\[ x^{2}+y^{2}+y=x^{2}+x^{2}\tan^{2}\left(x+\frac{\pi}{6}\right)+x\tan\left(x+\frac{\pi}{6}\right). \]
利用恒等式 \(\sec^{2}\theta=1+\tan^{2}\theta\)，有
\[ x\cdot x\sec^{2}\left(x+\frac{\pi}{6}\right)=x^{2}\left(1+\tan^{2}\left(x+\frac{\pi}{6}\right)\right)=x^{2}+x^{2}\tan^{2}\left(x+\frac{\pi}{6}\right). \]
所以左边为
\[ x\tan\left(x+\frac{\pi}{6}\right)+x^{2}+x^{2}\tan^{2}\left(x+\frac{\pi}{6}\right), \]
与右边相等，因此是解。

定义域说明：函数在 \(\cos\left(x+\frac{\pi}{6}\right)=0\) 的点处无定义，应在定义域内考虑。

答案：在定义域内是解

2. 题目：\(xy'=2y+1,\ y=5x^{2}+x\)

解答：验证函数是否满足微分方程。

计算导数：
\[ y'=10x+1. \]
代入左边：
\[ xy'=x(10x+1)=10x^{2}+x. \]
右边：
\[ 2y+1=2(5x^{2}+x)+1=10x^{2}+2x+1. \]
左边 \(10x^{2}+x\ne\) 右边 \(10x^{2}+2x+1\)，因此不是解。

答案：不是解

3. 题目：\(y''-\frac{2}{x}y'+\frac{2y}{x^{2}}=0,\ y=C_{1}x+C_{2}x^{2}\)

解答：验证函数是否满足微分方程。

计算导数：
\[ y'=C_{1}+2C_{2}x,\qquad y''=2C_{2}. \]
代入方程：
\[ y''-\frac{2}{x}y'+\frac{2y}{x^{2}}=2C_{2}-\frac{2}{x}(C_{1}+2C_{2}x)+\frac{2}{x^{2}}(C_{1}x+C_{2}x^{2}) \]
\[ =2C_{2}-\frac{2C_{1}}{x}-4C_{2}+\frac{2C_{1}}{x}+2C_{2}=0. \]
恒等于 \(0\)，因此是解。

定义域说明：方程在 \(x=0\) 处有奇点，解应在 \(x\ne0\) 的区间内成立。

答案：在 \(x\ne0\) 的区间内是解

4. 题目：\(xy'=y\left(1+\ln\frac{y}{x}\right),\ y=x\)

解答：验证函数是否满足微分方程。

计算导数：
\[ y'=1. \]
代入左边：
\[ xy'=x\cdot1=x. \]
右边：
\[ y\left(1+\ln\frac{y}{x}\right)=x\left(1+\ln\frac{x}{x}\right)=x(1+\ln1)=x(1+0)=x. \]
左边等于右边，因此是解。

定义域说明：对数函数要求 \(\frac{y}{x}>0\)，当 \(y=x\) 且 \(x>0\) 时成立。

答案：在 \(x>0\) 的区间内是解

3. 验证 \(y=C_{1}\cos x+C_{2}\sin x\)（\(C_{1},C_{2}\) 为任意常数）是方程 \(y''+y=0\) 的通解，并求满足初始条件 \(y|_{x=0}=1,\ y'|_{x=0}=3\) 的特解

解答：

首先验证是通解。计算导数：
\[ y'=-C_{1}\sin x+C_{2}\cos x,\qquad y''=-C_{1}\cos x-C_{2}\sin x. \]
代入方程：
\[ y''+y=(-C_{1}\cos x-C_{2}\sin x)+(C_{1}\cos x+C_{2}\sin x)=0. \]
恒成立，且解中含有两个任意常数，故是通解。

求特解。利用初始条件：
\[ y|_{x=0}=C_{1}\cos0+C_{2}\sin0=C_{1}=1\Rightarrow C_{1}=1, \]
\[ y'|_{x=0}=-C_{1}\sin0+C_{2}\cos0=C_{2}=3\Rightarrow C_{2}=3. \]
所以特解为
\[ y=\cos x+3\sin x. \]
答案：通解验证通过，特解为 \(y=\cos x+3\sin x\)

4. 写出下列曲线所满足的方程

1. 题目：曲线在点 \((x,y)\) 处的切线斜率等于该点横坐标的平方。

解答：切线斜率即导数 \(y'\)，所以
\[ y'=x^{2}. \]
答案：\(\frac{dy}{dx}=x^{2}\)

2. 题目：曲线上点 \(P(x,y)\) 处的法线与 \(x\) 轴的交点为 \(Q\)，且线段 \(PQ\) 被 \(y\) 轴平分。

解答：法线的斜率是切线斜率的负倒数。设切线斜率为 \(y'\)，则法线斜率为 \(-\frac{1}{y'}\)。

法线方程：
\[ Y-y=-\frac{1}{y'}(X-x). \]
与 \(x\) 轴交点 \(Q\)：令 \(Y=0\)，得
\[ 0-y=-\frac{1}{y'}(X_{Q}-x)\Rightarrow X_{Q}=x+yy'. \]
点 \(P\) 坐标为 \((x,y)\)，点 \(Q\) 坐标为 \((x+yy',0)\)。

线段 \(PQ\) 被 \(y\) 轴平分，即 \(PQ\) 的中点在 \(y\) 轴上。中点坐标为
\[ \left(\frac{x+(x+yy')}{2},\ \frac{y+0}{2}\right)=\left(x+\frac{yy'}{2},\ \frac{y}{2}\right). \]
中点在 \(y\) 轴上，故横坐标为 \(0\)：
\[ x+\frac{yy'}{2}=0\Rightarrow2x+yy'=0. \]
定义域说明：假设 \(y'\ne0\)（若 \(y'=0\) 需单独考虑）。

答案：\(2x+y\frac{dy}{dx}=0\)

3. 题目：曲线在点 \(P(x,y)\) 处的切线与 \(x\) 轴、\(y\) 轴的交点依次为 \(M\) 与 \(N\)，线段 \(MN\) 被点 \(P\) 平分，且曲线经过点 \((1,2)\)。

解答：切线方程：
\[ Y-y=y'(X-x). \]
与 \(x\) 轴交点 \(M\)：令 \(Y=0\)，得
\[ 0-y=y'(X_{M}-x)\Rightarrow X_{M}=x-\frac{y}{y'}. \]
与 \(y\) 轴交点 \(N\)：令 \(X=0\)，得
\[ Y_{N}-y=y'(0-x)\Rightarrow Y_{N}=y-xy'. \]
点 \(P\) 是线段 \(MN\) 的中点，故
\[ x=\frac{X_{M}+0}{2}=\frac{x-\frac{y}{y'}}{2},\qquad y=\frac{0+Y_{N}}{2}=\frac{y-xy'}{2}. \]
由第一个方程：
\[ 2x=x-\frac{y}{y'}\Rightarrow x=-\frac{y}{y'}\Rightarrow xy'+y=0. \]
由第二个方程：
\[ 2y=y-xy'\Rightarrow y=-xy'\Rightarrow xy'+y=0. \]
两者等价，所以微分方程为
\[ x\frac{dy}{dx}+y=0. \]
此外，曲线经过点 \((1,2)\)，这是一个初始条件。

求解显式解：
\[ x\frac{dy}{dx}+y=0\Rightarrow\frac{dy}{dx}=-\frac{y}{x}. \]
分离变量：
\[ \frac{dy}{y}=-\frac{dx}{x}. \]
积分：
\[ \ln|y|=-\ln|x|+C\Rightarrow\ln|y|=\ln\left|\frac{1}{x}\right|+C\Rightarrow y=\frac{K}{x}, \]
其中 \(K=\pm e^{C}\) 为常数。

代入初始条件 \(y(1)=2\)：
\[ 2=\frac{K}{1}\Rightarrow K=2. \]
所以特解为
\[ y=\frac{2}{x}. \]
定义域说明：解在 \(x\ne0\) 的区间内成立。

答案：微分方程为 \(x\frac{dy}{dx}+y=0\)，满足 \(y(1)=2\) 的特解为 \(y=\frac{2}{x}\)
